\chapter*{Preface}

In science fiction \textit{quantum} is often used to explain away a technology that seems to defy the laws of physics. In the 2018 superhero movie \textit{Ant-Man}, the protagonist is able to become smaller and smaller until they enter the mysterious \textit{Quantum Realm}, described as “a reality where all concepts of time and space become irrelevant”. Perhaps the writers have a consistent set of logical explanations to describe how their worlds work. Our notions of space and time in our world might not be true in their world. 

The set of rules that define quantum mechanics are a subject of intense creative destruction. Theoretically, defining space and time on the quantum level is a subject for higher dimensions and exotic symmetries. Experimentally, pushing the limits of technology to define every greater precision in measurement has given birth to the current information revolution. This revolution, like the industrial revolution, or the digital revolution that brought forth the world wide web, is as exciting as it is poorly understood. The quantum revolution offers exciting new technologies that could transform how we view the world (sensing), process it (computing) and share those insights (communications). 

No single book could do justice to all the categories of quantum technologies, so this book will focus on computing in the most concise explanation, accepting the simple reality that the field is changing too rapidly for most of the content to be anywhere close to state-of-the-art. If all this book can accomplish for the reader is to develop an honest and intuitive appreciation for quantum information processing, it would have been worth writing. Least of all for myself. 
